\begin{table}[H]
\centering
\caption{\textbf{More valid labels and more balanced substrate weights make separate contributions.} In muscle development data, recovering eligible high/low labels improves prediction under row-wise sampling. Pair-first weighting adds a further gain with the label pool and class balance fixed. The same-fold comparisons isolate the label and weighting effects; intervals are pair-clustered bootstrap intervals for each indicated step, in percentage points.}\label{tab:supervisiondecomposition}
\begin{tabular}{lrrr@{\hspace{2\tabcolsep}\hspace{0.6em}}rr}
\toprule
Metric & Original reference & All high/low, row-wise & $\Delta$ from labels [95\% CI] & All high/low, pair-first & $\Delta$ from pair weighting [95\% CI] \\
\midrule
AUROC & 0.8456 & 0.8659 & +2.04 [+1.06, +3.00] & 0.8885 & +2.27 [+1.25, +3.36] \\
AUPRC & 0.8256 & 0.8520 & +2.63 [+1.25, +4.17] & 0.8730 & +2.08 [+0.76, +3.55] \\
F1 & 0.7783 & 0.7924 & +1.42 [+0.21, +2.57] & 0.8208 & +2.85 [+1.68, +4.01] \\
\bottomrule
\end{tabular}
\end{table}
